\section{系统总体设计}

在这一章中，我们从整体上介绍系统的设计理念、整体框架、功能模块组成以及实验的完整流程。

\subsection{总体设计思路}

Thought2Text 数据集共有 40 个类别，大约包含 12000 条 trial，这个规模大约只有主流视觉语言模型预训练集的百万分之一。数据量这么小，直接从 EEG 信号学出生成自然文本的能力是很困难的。另外，这个数据集也没有人工写的图片描述，只有类别标签这一种原始标注。所以我们把整个系统设计成四个递进的阶段，从简单到复杂，从定量验证逐步过渡到生成式输出。每个阶段我们都有明确的输入输出定义，中间结果也是可验证的，前面阶段的输出可以直接用做后面阶段的输入或者基础模块。

\textbf{第一阶段：CLIP 语义空间构建}。我们用 CLIP 在大规模数据上学到的图文对齐空间来做跨模态的桥梁。40 个类别的文本 prompt 先通过 CLIP text encoder 算好，存成固定的文本原型向量；所有图片也提前用 CLIP ViT 提取好视觉特征。这样一来，后面所有的 EEG 和图像之间的语义交互就都有了统一的操作空间，不需要在少量 EEG 数据上从头去学语义结构了。

\textbf{第二阶段：A2 EEG 语义融合}。我们设计了一个 A2 temporal-spectral-spatial EEG encoder，从 EEG 信号里分别提取时间、频谱和空间三个维度的语义信息。然后在 CLIP 语义空间中，用门控残差融合（gated residual fusion）把 EEG embedding 和经过退化处理的图像视觉 embedding 结合起来，主要用分类损失来训练。我们还设置了 shuffled EEG 和 random EEG 两种对照实验，用来严格验证 EEG 里面确实携带了真实的语义信息。

\textbf{第三阶段：token 级别的 EEG 视觉增强}。第二阶段做的融合是在 pooled embedding 层面——整张图片的视觉信息被压缩成了一个向量。为了让 EEG 信息能够进入生成式 VLM 的 token 处理流程，我们在这个阶段设计了 VTF 模块，直接在 CLIP ViT 的 visual patch token 层注入 EEG token，通过 cross-attention 机制生成 EEG-enhanced vision tokens。

\textbf{第四阶段：生成式 VLM 集成与验证}。我们把 enhanced vision tokens 送进 Qwen2-VL-2B-Instruct 生成式模型，用 LoRA 做参数高效的微调。我们对比了好几种生成方案和训练目标，再用 best-of-N reranking 来提升输出质量，最终实现了从退化图像和 EEG 信号到自然语言描述的完整通路。

\subsection{系统总体架构}

系统的整体结构如图~\ref{fig:overall-system} 所示。

\placeholderfigure{系统总体架构图}{overall-system-placeholder}{4.5cm}

数据在系统中的流向是这样的：退化图像经过冻结的 CLIP ViT 编码器处理之后，输出一个 512 维的 pooled image embedding 和一组 $50\times512$ 维的 visual patch tokens；EEG trial 经过 A2 encoder 处理之后，输出一个 512 维的 pooled EEG embedding 和一组 $4\times512$ 维的 EEG tokens。在语义预测分支里，两个 pooled embedding 经过 gated fusion 融合，再跟文本原型做相似度匹配，得到类别预测结果。在 token 级别的分支里，visual tokens 和 EEG tokens 通过 VTF cross-attention 生成 enhanced vision tokens。在生成式分支里，增强后的视觉 token 先经过我们实现的轻量 prefix adapter 映射到 Qwen2-VL 能接收的嵌入空间，然后跟语义提示文本一起输入生成式解码器，经过 LoRA 微调之后输出自然语言描述。

\subsection{模块划分}

整个系统由七个核心模块组成，各个模块的功能说明列在表~\ref{tab:modules} 中。

\begin{table}[htbp]
  \centering
  \caption{系统模块划分及功能说明}
  \label{tab:modules}
  \begin{tabularx}{\linewidth}{l X}
    \toprule
    模块名称 & 功能说明 \\
    \midrule
    图像编码模块 & CLIP ViT-B/32，输出 pooled embedding 和 visual patch tokens，模型参数固定不更新 \\
    EEG 编码模块 & A2 temporal-spectral-spatial encoder，输出 pooled EEG embedding 和 EEG tokens \\
    文本原型模块 & 用 CLIP text encoder 预先算好 40 类文本原型向量，训练阶段保持不变 \\
    语义融合模块 & 用 Gated residual fusion 把图像和 EEG 的 embedding 融合起来，再跟文本原型匹配做分类 \\
    Token 融合模块 & 用 VTF cross-attention 把 EEG tokens 融合进 visual tokens，生成增强后的视觉 token \\
    生成式描述模块 & Qwen2-VL-2B-Instruct + LoRA，接收增强视觉 token 并生成自由格式的文本描述 \\
    评估模块 & 多种模式对照评测，统计指标，找出有代表性的定性案例 \\
    \bottomrule
  \end{tabularx}
\end{table}

\subsection{实验流程}

图~\ref{fig:experiment-flow} 展示了从数据准备到最后评估的完整实验流程。

\placeholderfigure{实验流程图}{experiment-flow-placeholder}{4.0cm}

整个流程包括以下几个步骤：(1) 数据准备——下载 EEG 数据、关联 ImageNet 图片、构建 JSONL 格式的 manifest 文件；(2) CLIP 特征预缓存——提前算好文本原型和所有图片的视觉特征；(3) A2 EEG encoder 训练——用 InfoNCE 对比损失函数做 EEG 到 CLIP 的对齐训练；(4) 语义融合训练——把 gated fusion 和分类模块放在一起联合训练，在六种不同的退化条件下做多模式评估；(5) VTF token 融合训练——训练 cross-attention 融合模块；(6) 生成式模型训练——把 enhanced vision tokens 输入 Qwen2-VL，用 LoRA 做微调，比较不同的训练目标和 LoRA 配置；(7) 全量测试与 reranking——在测试集上把所有组合都跑一遍评估，对效果最好的 checkpoint 做 best-of-N reranking；(8) 报告生成——汇总各项指标，挖掘有代表性的定性案例。

所有实验都在单张 GPU（NVIDIA RTX 6000 Ada，48GB 显存）上运行，使用 bfloat16 混合精度训练，没有用多 GPU 分布式训练或者 4-bit 量化。
